\section{An Architecture for a Blockchain-Based Self Protecting System}\label{sec:architecture}

\begin{figure}
\centering
\includegraphics[width=0.7\linewidth]{figures/irs-blockchain-architecture.pdf}
\caption{Reference architecture from~\cite{albero2026blockchain}, see description in the text.}\label{fig:irs-block-architecture}
\end{figure}

Our reference architecture is based on the one proposed in~\cite{albero2026blockchain}, which follows the classic autonomic computing paradigm of Monitor, Analyze, Plan, and Execute (MAPE) composed of a managed system and a manager system~\cite{kephart2003vision, yuan2014systematic}. Both systems consist of sets of containerized components deployed on physical or virtual hosts.
A comprehensive view of the architecture is shown in Figure~\ref{fig:architecture}.

The managed system comprises application containers whose behavior is characterized by monitored state variables. The set of these variables defines the monitored system state \circled{2}~\cite{iannucci2018model} that is affected by some security event \circled{1}.

The manager system comprises IDSs \circled{3}, actuators \circled{8}, and blockchain nodes \circled{4} and \circled{7}, each deployed in a separate container. In the original paper~\cite{albero2026blockchain}, blockchain nodes are Quorum nodes~\cite{Quorum56:online}. Together, the blockchain and the actuators constitute the Intrusion Response System (IRS). Each IDS and each actuator is associated with a blockchain node that enables interaction with the blockchain, while additional blockchain nodes may be deployed to increase resilience.

The core idea is that the blockchain, through a smart contract, maintains a shared state \circled{5} across IDSs and actuators. This shared state includes a representation of the monitored system state, as well as information about the internal state of the IRS. More precisely, the smart contract implements finite-state-machine logic and must be realized correctly and securely to guarantee integrity~\cite{atzei2017survey, kushwaha2022systematic}.

IDSs submit transactions reporting their view of one or more monitored variables. Because IDSs themselves may be compromised, the smart contract updates the shared state only through a consensus-based voting mechanism, enabling the system to tolerate the compromise of up to one-third of the sensors~\cite{wang2022bft}. For this reason, each variable should ideally be monitored by at least three heterogeneous IDSs. Whenever the shared state is updated, the smart contract applies the response logic \circled{6}, namely a state–action map used to determine the appropriate countermeasure.

Accepted transactions are broadcast to all blockchain nodes, and state changes trigger countermeasures through push-based mechanisms \cite{buterin2013ethereum}. Finally, for security reasons, actuators are assumed to be blind: they execute their tasks without reading any state directly from the managed system, thereby preventing direct attacks from compromised application containers.

\subsection{Related Works}
The literature on Self-Protecting Systems has been systematically reviewed by Yuan et al. \cite{yuan2014systematic}, who provide a taxonomy of detection and response strategies, and by Tziakouris et al. \cite{tziakouris2018survey}, who broaden the scope to self-adaptive security in large-scale open environments. Both works highlight that securing the autonomic manager itself remains an open challenge, which motivates the blockchain-based approach adopted in \cite{albero2026blockchain} and revisited here. The use of distributed ledger technologies for intrusion detection has been extensively surveyed by Li et al. \cite{li2022surveying}, who classify collaborative IDS architectures and discuss how blockchain-based trust management mitigates insider threats such as Sybil and collusion attacks. Concrete instantiations of this idea, such as the deep blockchain framework proposed by Alkadi et al. \cite{alkadi2021deep}, secure the \emph{detection plane} by replicating detection logic across a permissioned ledger, in contrast, our work focuses on the \emph{management plane}, where the autonomic manager coordinates detection and response.